\chapter{Introduction}

\section{What is Giovi}

Giovi (pronounced as \textit{"jovi"}) is a high level, multipurpose,
strongly typed, imperative programming and scripting language.

\begin{listing}
	\label{lst:hello_world}
	\begin{minted}{\giovilexer}
"hello world" -> _out
	\end{minted}
	\caption{The usual sample program that prints "hello world" to stdout}
\end{listing}

Giovi is originally born as a configuration language and as such, it relies on
a mostly structural type system; it puts a lot of emphasis on the use of few
basic data structures that, combined with simple yet powerful syntactic
features, allows writing code that is easy to read and fast to write.

\begin{listing}
	\label{lst:intro_overview}
	\caption[Syntax intro overview]{
		A variable declaration of type \code{User}, initialized with a variety
		of data types: primitive types (numbers, strings, boolean), collection
		types (lists, maps), product types and functions
	}
	\begin{minted}{\giovilexer}
var simonvic <- User.<
	id: 69
	username: "simonvic"
	aliaes: ["Simone", "Simon"]
	bio: "{
		hello I'm simonvic,
		and I'm working on Giovi!
		}"
	age: 26
	favoriteColor: Rgb.<0, 255, 0>
	credits: 69.420
	roles: [
		"admin" => true
		"devel" => true
		"guest" => false
	]
	onLogin: function(user User) {
		"hello ${user.username}" -> _out
		if "admin" in user.roles do {
			"Granting admin priviledges" -> _out
		}
	}
>
	\end{minted}
\end{listing}

Giovi evolved to be more than a configuration language, by providing
user-defined data structures, first-class functions and types. Giovi wants to
experiment with the idea of non-nullability, where everything has a
(overwritable) "default" value that, along the support for automatic data
widening (but explicit narrowing), allows writing flexible code with few
keystrokes.

\begin{listing}
	\label{lst:calculator}
	\caption{
		An example of a \code{calculator} function and the use of functions as
		first-class citizens
	}
	\begin{minted}{\giovilexer}
plus(x Real, y Real) => Real { return x + y }
minus(x Real, y Real) => Real { return x - y }
times(x Real, y Real) => Real { return x * y }
divided(x Real, y Real) => Real { return x / y }

#? Calculator function
#? leftHandSide: first operand
#? strategy: a function with signature (Real,Real)=>Real
#? rightHandSide: second operand
calculator(
	leftHandSide Real,
	strategy function(x Real, y Real) => Real,
	rightHandSide Real,
) => Real {
	# apply strategy
	return strategy(leftHandSide, rightHandSide)
}

calulator(5, plus, 2) -> _out # 5 + 2
calulator(5, minus, 2) -> _out # 5 - 2
calulator(5, times, 2) -> _out # 5 * 2
calulator(5, divided, 2) -> _out # 5 / 2
	\end{minted}
\end{listing}

The grammar design philosophy is to have a language that is easy to read and
fast to write. The syntax aims to have minimal lexical and syntactic clutter
where keywords are preferred over symbols, and to provide many quality-of-life
features that speed up the writing process, without hindering the reading by
relying on the use of symbols that is consistent and predictable.

\section{Executing Giovi code}

Giovi code can be executed from a file

\begin{minted}{Bash}
giovi some_script.giovi
\end{minted}

Alternatively, the source code can be passed directly to the CLI

\begin{minted}{Bash}
giovi -f '"hello world" -> _out'
\end{minted}

\subsection{REPL}

Giovi also has a minimal REPL (Read Evaluation Print Loop), that allows you to
write and execute code on the fly

\begin{minted}{Bash}
giovi --repl
> "The sum of 10 and 20 is equal to ${10 + 20}" -> _out
> The sum of 10 and 20 is equal to 30
\end{minted}

